\subsection{LINEAR EQUATIONS OF ORDER n WITH CONSTANT COEFFICIENTS}

If the coefficients \(a_{j}\) in equation (33) are constant, the corresponding matrix \(A\) is constant; the characteristic equation is obtained by writing that \(e^{rt}\) is a solution, which gives

\[
r ^ {n} - a _ {1} r ^ {n - 1} - \ldots - a _ {n - 1} r - a _ {n} = 0.\tag{40}
\]

Let \(r_j\) ( \(1 \leqslant j \leqslant q\) ) be the distinct roots of this equation, and \(n_j\) ( \(1 \leqslant j \leqslant q\) ) the multiplicity of the root \(r_j \left( \sum_{j=1}^{q} n_j = n \right)\) . By the results of IV, p.~188 to 194, to each root \(r_j\) there corresponds, for the homogeneous equation

\[
\mathrm{D} ^ {n} x - a _ {1} \mathrm{D} ^ {n - 1} x - \dots - a _ {n - 1} \mathrm{D} x - a _ {n} x = 0\tag{41}
\]

a system of \(n_j\) linearly independent integrals

\[
u _ {j k} (t) = e ^ {r _ {j} t} p _ {j k} (t),
\]

where \(p_{jk}\) is a polynomial (with complex coefficients) of degree \(\leqslant n_{j}-1\) ( \(1\leqslant k\leqslant n_{j}\) ); further, the n integrals \(u_{jk}\) ( \(1\leqslant j\leqslant q\) , \(1\leqslant k\leqslant n_{j}\) ) so obtained are linearly independent. It follows that the \(n_{j}\) polynomials \(p_{jk}\) ( \(1\leqslant k\leqslant n_{j}\) ) are linearly independent in the space of polynomials in t of degree \(\leqslant n_{j}-1\) , so form a basis (over C) of this space, since the latter is of dimension \(n_{j}\) . In other words:

PROPOSITION 9. Let \(r_j\) ( \(1 \leqslant j \leqslant q\) ) be the distinct roots of the characteristic equation (40), and let \(n_j\) be the multiplicity of the root \(r_j\) ( \(1 \leqslant j \leqslant q\) ). Then the \(n\) functions \(t^k e^{r_j t}\) ( \(1 \leqslant k \leqslant n_j\) , \(1 \leqslant j \leqslant q\) ) are linearly independent integrals of the homogeneous equation (41).

One can prove this result directly in the following way. It follows from equation (41) that the \(n^{th}\) derivative of every integral of this equation is differentiable on R, from which one deduces immediately, by induction on the integer m \textgreater{} n, that every integral of (41) admits a derivative of order m, that is, is indefinitely differentiable on R. Let D be the (non-topological) vector space over C of indefinitely differentiable complex-valued functions on R; the map \(x \mapsto Dx\) is an endomorphism of this space, and equation (41) can be written

\[
f (\mathrm{D}) x = 0\tag{42}
\]

\[
\text { where } f (\mathrm{D}) = \mathrm{D} ^ {n} - a _ {1} \mathrm{D} ^ {n - 1} - \dots - a _ {n - 1} \mathrm{D} - a _ {n} (\mathrm{Alg}., \mathrm{IV}, \mathrm{p}. 4).
\]

PROPOSITION 10. Let \(g\) and \(h\) be two relatively prime polynomials such that \(f = gh\) . The subspace of solutions of (42) is the direct sum of the subspaces of solutions of the two equations

\[
g (\mathrm{D}) x = 0, \quad h (\mathrm{D}) x = 0.
\]

Now, by the Bezout identity (Alg., VII. 2, th. 1) there exist two polynomials \(p(\mathrm{D})\) and \(q(\mathrm{D})\) such that \(p(\mathrm{D})g(\mathrm{D}) + q(\mathrm{D})h(\mathrm{D}) = 1\) . For every solution \(x\) of (42) one can write \(x = y + z\) , where \(y = p(\mathrm{D})g(\mathrm{D})x\) and \(z = q(\mathrm{D})h(\mathrm{D})x\) ; and then \(h(\mathrm{D})y = p(\mathrm{D})(f(\mathrm{D})x) = 0\) and \(g(\mathrm{D})z = q(\mathrm{D})(f(\mathrm{D})x) = 0\) . On the other hand, if both \(g(\mathrm{D})x = 0\) and \(h(\mathrm{D})x = 0\) , one deduces that

\[
x = p (\mathrm{D}) (g (\mathrm{D}) x) + q (\mathrm{D}) (h (\mathrm{D}) x) = 0,
\]

which completes the proof.

With the preceding notation one can write

\[
f (\mathrm{D}) = \prod_ {j = 1} ^ {q} (\mathrm{D} - r _ {j}) ^ {n _ {j}}
\]

and prop. 10, by induction on \(q\) , shows that the subspace of solutions of (42) is the direct sum of the subspaces of solutions of the \(q\) equations

\[
(\mathrm{D} - r _ {j}) ^ {n _ {j}} x = 0 \quad (1 \leqslant j \leqslant q).\tag{43}
\]

Now, for every complex number r one has

\[
\mathrm{D} (e ^ {r t} x) = e ^ {r t} (\mathrm{D} + r) x\tag{44}
\]

so (43) is equivalent to

\[
\mathrm{D} ^ {n _ {j}} \left(e ^ {- r _ {j} t} x\right) = 0
\]

and thus has as solutions the functions \(e^{r_{j}t} p_{j}(t)\) , where \(p_{j}\) runs through the set of polynomials of degree \(\leqslant n_{j}-1\) ; thus one recovers prop. 9 of IV, p.~194.

Assuming that the homogeneous equation (41) has been solved (that is, assuming that the characteristic roots have been found), one knows that the method of variation of parameters allows one to find the solutions of the inhomogeneous equation

\[
\mathrm{D} ^ {n} x - a _ {1} \mathrm{D} ^ {n - 1} x - \dots - a _ {n - 1} \mathrm{D} x - a _ {n} x = b (t)\tag{45}
\]

where \(b(t)\) is an arbitrary regulated function (IV, p.~182); we note that if \(b(t)\) is indefinitely differentiable on an interval \(J\) then all the integrals of (45) are indefinitely differentiable on J. In the particular case \(b(t) = e^{\alpha t} p(t)\) , where p is a polynomial (with complex coefficients) and \(\alpha\) is an arbitrary complex number, one obtains an integral of (45) more simply by the following method. Put \(x = e^{\alpha t} y\) ; the equation

\[
f (\mathrm{D}) x = e ^ {\alpha t} p (t)
\]

can, by (44), be written

\[
f (\alpha + \mathrm{D}) y = p (t)
\]

or, by Taylor's formula applied to the polynomial \(f(\mathbf{D})\) ,

\[
\frac {f ^ {(n)} (\alpha)}{n !} \mathrm{D} ^ {n} y + \frac {f ^ {(n - 1)} (\alpha)}{(n - 1) !} \mathrm{D} ^ {n - 1} y + \dots + \frac {f ^ {\prime} (\alpha)}{1 !} \mathrm{D} y + f (\alpha) y = p (t).\tag{46}
\]

Let \(m\) be the degree of the polynomial \(p(t) = \sum_{k=0}^{m} \lambda_k t^{m-k}\) ; if \(f(\alpha) \neq 0\) (that is, if \(\alpha\) is not a characteristic root), there exists one and only one polynomial \(u(t) = \sum_{k=0}^{m} c_k t^{m-k}\) of degree \(m\) which is a solution of (46), for the coefficients \(c_k\) are determined by the system of linear equations

\[
\begin{array}{r l} f (\alpha) c _ {k} + \binom {m - k + 1} {1} f ^ {\prime} (\alpha) c _ {k - 1} + \binom {m - k + 2} {2} f ^ {\prime \prime} (\alpha) c _ {k - 2} + \dots \\ & + \binom {m} {k} f ^ {(k)} (\alpha) c _ {0} = \lambda_ {k} \quad (0 \leqslant k \leqslant m) \end{array}
\]

which clearly admits one and only one solution. If, on the other hand, \(\alpha\) is a characteristic root, and h is its multiplicity, the preceding calculation shows that there is one and only one polynomial of degree m such that every solution of \(D^{h}y = v(t)\) is an integral; in other words, every polynomial solution of (46) is then of degree \(m + h\) (``resonance'').

